\section{Submission Overview}
\label{sec:submission-overview}

This revised draft is organized as the first-stage NGCC submission package for \emph{Kylin}, the submission name used in this repository for the \SPHINCSA-based construction. It preserves the original technical specification of \SPHINCSA\ while making the submission-oriented items explicit. Whenever the current repository does not yet contain finalized evidence or deliverables, we mark the gap with \todo{...}.

\subsection{Algorithm Basic Information}

\begin{table}[htbp]
    \centering
    \caption{Basic information for the current Kylin submission package.}
    \label{tab:basic-info}
    \begin{tabular}{p{0.30\linewidth}p{0.62\linewidth}}
    \toprule
    Item & Current draft status \\ \midrule
    Algorithm name & Kylin \\
    Algorithm function & Digital signature \\
    Algorithm family & Stateless hash-based signature \\
    Core algorithms & Key generation, signature generation, signature verification \\
    Baseline construction & Kylin is the NGCC submission name for the \SPHINCSA\ construction, itself derived from \SPHINCSP\ with \WOTSCS\ replacing \WOTSP\ and with retuned parameter sets \\
    Hash/XOF interface in the packaged code & ICCS auxiliary functions, with domain-separated wrappers built from \texttt{pseudoXOF}; the legacy draft discussion still records the SHAKE- and SHA2-based research instantiations \\
    Parameter sets currently included & Submitted NGCC instances \texttt{Kylin-128}, \texttt{Kylin-256}, and \texttt{Kylin-512}; the legacy \texttt{sphincs-a-*} families remain in the draft as historical baselines; see \cref{tab:parameters,tab:ngcc-parameters} \\
    Submitters & \todo{List the names, affiliations, postal addresses, phone numbers, email addresses, and signatures of all submitters.} \\
    Primary contact & \todo{Identify the submission contact and provide the signed contact information required by the call.} \\
    \bottomrule
    \end{tabular}
\end{table}

\subsection{Conformance Notes for the Current Draft}

\begin{itemize}
    \item The required algorithm description is provided by the technical sections on notation, the one-time signature layer, the hypertree, FORS, the full \SPHINCSA\ construction, and the parameter/hash-function instantiations; see \cref{sec:notation,sec:ots,sec:hypertree,sec:fors,sec:spx-core,sec:instantiations}.
    \item The required design rationale, security analysis, performance evaluation, and feature statement are provided in \cref{sec:wots-optimization,sec:design-rationale,sec:security-claims,sec:performance,sec:feature-statement}.
    \item The legacy parameter tables inherited from the original \SPHINCSA\ draft are still useful technical baselines, but this revised submission draft now additionally includes NGCC-oriented families covering the required 128/256/512-bit classical levels with corresponding 80/128/256-bit generic quantum levels; see \cref{tab:ngcc-parameters}. The selection is reproducible from \texttt{paper/code/search\_ngcc\_params.py}.
    \item The recommended NGCC parameter families are tuned with $q=2^{64}$ in the concrete-security formulas and all use hypertree height $h=70$, so the submission text can now state support for at least $2^{64}$ signatures under one key pair in the intended single-key usage model.
    \item The NGCC SIG wrapper in the packaged code accepts message lengths up to $2^{60}$ bytes, i.e.\ $2^{63}$ bits, which matches the bound stated in the call.
    \item Separate submission artifacts required by the call, including signed forms and ownership/IP statements, are tracked in Appendix~\ref{apx:submission-checklist}. The packaged ISO C reference implementation, packaged optimized implementation, and generated test vectors are now present in the repository.
\end{itemize}

\subsection{Legacy Draft Families vs. Recommended NGCC Families}

The current repository still contains the original \SPHINCSA\ draft parameter families. These families remain useful historical baselines for studying size/speed trade-offs, but they are not identical to the target families requested by the NGCC call. In particular, the legacy tables are organized around approximately 128/192/255-bit classical security, whereas the call asks for explicit 128/256/512-bit classical security families with corresponding quantum-security targets.

Using the generic formulas currently stated in the draft, one can infer an approximate generic quantum bit-security by halving the corresponding classical bit-security. This is only an inference from the present formulas, not yet a finalized NGCC security claim.

\begin{table}[htbp]
    \centering
    \caption{Alignment of the current draft parameter families with the NGCC targets. Each row applies to both SHAKE and SHA2 instantiations.}
    \label{tab:ngcc-alignment}
    \begin{tabularx}{\linewidth}{l c c c c X}
    \toprule
    Family & Classical bitsec & Approx.\ quantum bitsec & $h$ & Sig.\ bytes & Current alignment status \\ \midrule
    128s & 128 & $\approx 64$ & 63 & 6880 & \todo{Not yet a finalized NGCC 128/80 claim; the current $h=63$ setting also makes the required $2^{64}$-message support statement non-obvious.} \\
    128f & 128 & $\approx 64$ & 63 & 16720 & \todo{Not yet a finalized NGCC 128/80 claim; the quantum-security target still needs explicit NGCC-oriented justification.} \\
    192s & 192 & $\approx 96$ & 63 & 14568 & \todo{Intermediate NIST-style family only; not itself an NGCC target family.} \\
    192f & 192 & $\approx 96$ & 64 & 34896 & \todo{Intermediate NIST-style family only; not itself an NGCC target family.} \\
    256s & 255 & $\approx 127$ & 66 & 27232 & \todo{Close to a 256-bit classical family, but still not a finalized NGCC 256/128 claim and still needs explicit submission wording.} \\
    256f & 255 & $\approx 127$ & 68 & 49312 & \todo{Close to a 256-bit classical family, but still not a finalized NGCC 256/128 claim and still needs explicit submission wording.} \\
    \bottomrule
    \end{tabularx}
\end{table}

For the actual NGCC submission package, this revised draft now uses a first-stage three-family Kylin parameter set generated by the repository-local parameter script. The script enumerates feasible candidates using the formulas in \cref{sec:security-claims}, computes exact \WOTSCS lengths, sets $q=2^{64}$ to match the call's per-key signing-capacity requirement, and emits a fixed common implementation shape while the final parameter-selection method remains under discussion. The resulting common shape is
\[
    (h,d,a,w) = (70,10,14,64).
\]

\begin{table}[htbp]
    \centering
    \caption{First-stage NGCC-aligned Kylin parameter families emitted by \texttt{paper/code/search\_ngcc\_params.py}. The 128/80 row records the current design decision to use $n=16$; its detailed quantum-security justification remains part of the parameter discussion.}
    \label{tab:ngcc-parameters}
    \scriptsize
    \begin{tabular}{l c c c c c c c c c c c c c}
    \toprule
    Family & NGCC lvl & $n$ & $h$ & $d$ & $\log t$ & $k$ & $w$ & $l$ & bitsec & q-bitsec & pk & sk & sig \\ \midrule
    \csvreader[head to column names]{paper/data/para_ngcc.csv}{}%
    { \name & \ngccclassical/\ngccquantum & \n & \h & \d & \a & \k & \w & \len & \classical & \quantum & \pk & \sk & \sig \\ }
    \vspace{-0.8em} \\ \bottomrule
    \end{tabular}
\end{table}

The corresponding rough hash-call estimates from the search script are approximately $2.90 \times 10^6$, $5.06 \times 10^6$, and $9.26 \times 10^6$ for the recommended 128-, 256-, and 512-bit NGCC families, respectively.

The selected families were chosen for the following reasons:

\begin{itemize}
    \item The shared structure $(h,d,a,w)=(70,10,14,64)$ keeps the specification and implementation uniform across all three NGCC levels; only $n$, $k$, and the derived \WOTSCS length $l$ change.
    \item The common hypertree height $h=70$ is comfortably above the call's minimum $2^{64}$-message requirement and makes the signing-capacity statement straightforward in the single-key model used by the tuning script.
    \item The choice $w=64$ gives the current \WOTSCS lengths $l \in \{23,45,87\}$ for the \texttt{Kylin-128}, \texttt{Kylin-256}, and \texttt{Kylin-512} instances.
    \item The legacy 192-bit family remains useful for research comparison, but it is no longer part of the recommended NGCC submission family.
    \item \todo{If the final package wants to mirror the evaluation document's broader ``at most $2^{80}$ queries'' convention literally in the concrete-security tables, add a separate note clarifying the relation between that convention and the single-key $q=2^{64}$ tuning used here.}
\end{itemize}

\subsection{Correctness and Failure Probability}

\SPHINCSA\ is a stateless hash-based signature scheme. Under honest execution, the signing and verification algorithms are deterministic once the per-signature randomness is fixed, and the design does not introduce a decryption- or decapsulation-style failure event. In that sense, there is no intrinsic non-negligible correctness failure probability parameter analogous to code-based or lattice-based decapsulation failure.

Accordingly, for every supported instance, every key pair output by the packaged key-generation algorithm, and every message of length at most $2^{63}$ bits, verification under the matching public key accepts every signature produced by the matching signing algorithm except in the event of implementation faults or a failure of the underlying cryptographic assumptions. The current package therefore treats correctness as a deterministic algorithmic property of the implementation, with no separate decryption-style failure term to report.

\subsection{Status of External Deliverables}

The current repository already contains a detailed LaTeX specification, pseudocode listings, benchmark data, parameter-evaluation scripts, the NGCC implementation package, and generated test vectors. However, the call still requires several external materials in a stricter finalized form, especially signed forms, ownership/IP declarations, and further submission-facing performance evidence. A concise checklist of the remaining deliverables is provided in Appendix~\ref{apx:submission-checklist} so that the remaining work can be tracked directly inside the draft.
